% VI37-DETAILED-PROBLEM-03
\begin{problem}[VI.37.P3: Intersection multiplicity preview]\label{prob:vi37-03}
Inside \(\mathbb P^2\), intersect the line \(Y=0\) with the conic \(Y^2-XZ=0\). Determine the scheme-theoretic intersection.
\end{problem}

\ifdefined\FullProblemDossiers
\begin{solution}
Let
\[
I_L=(Y),
\qquad
I_C=(Y^2-XZ).
\]
Scheme-theoretic intersection is defined by the sum:
\[
I_L+I_C
=
(Y,Y^2-XZ)
=
(Y,XZ).
\]
Therefore
\[
L\cap C
=
\Proj k[X,Y,Z]/(Y,XZ).
\]
Inside the line \(Y=0\cong\mathbb P^1\), the equation is \(XZ=0\), whose homogeneous ideal decomposes as
\[
(XZ)=(X)\cap(Z).
\]
Thus
\[
\boxed{
L\cap C
=
\{[1:0:0],[0:0:1]\}
}
\]
as a reduced two-point scheme. On the \(X\)-chart the second point is absent and the local ideal is \((z)\); on the \(Z\)-chart the first is absent and the local ideal is \((x)\). Hence both points have length one.
\end{solution}
\fi
